\chapter{Nystagmus}

\section*{Definition}
Nystagmus is an involuntary, rhythmic oscillation of the eyes characterized by slow and fast phases (jerk nystagmus) or pendular oscillations of equal velocity. It may be congenital or acquired, and is classified by direction (horizontal, vertical, torsional), waveform (jerk versus pendular), and the effect of gaze position and fixation.

\section*{What Happens in Your Body}
Nystagmus arises from disruption of the neural integrator mechanisms that maintain stable gaze holding. The vestibular-ocular reflex (VOR), smooth pursuit, and neural integrator circuits in the brainstem and cerebellum work together to stabilize retinal images. Lesions affecting the vestibulocerebellum (flocculus, nodulus), medial vestibular nucleus, nucleus prepositus hypoglossi, or interstitial nucleus of Cajal impair gaze-holding ability, causing the eyes to drift (slow phase) followed by corrective saccades (fast phase). Congenital nystagmus involves instability of the oculomotor system without structural brainstem pathology.

\section*{Causes}
\begin{itemize}
  \item Peripheral vestibular causes: Benign paroxysmal positional vertigo (BPPV), vestibular neuritis, M\'eni\`ere disease, labyrinthitis, acoustic neuroma, ototoxicity
  \item Central nervous system causes: Brainstem stroke (lateral medullary syndrome), cerebellar infarction or hemorrhage, multiple sclerosis, Chiari malformation, brainstem tumors, Wernicke encephalopathy, paraneoplastic syndromes
  \item Congenital causes: Infantile nystagmus syndrome (idiopathic or associated with albinism, aniridia, achromatopsia, optic nerve hypoplasia, Leber congenital amaurosis)
  \item Toxic and metabolic: Anticonvulsants (phenytoin, carbamazepine), lithium, alcohol intoxication, barbiturates, carbon monoxide poisoning, thiamine deficiency
  \item Other causes: Severe visual loss (sensory nystagmus), spasmus nutans (children), see-saw nystagmus (suprasellar lesions)
\end{itemize}

\section*{When to See a Doctor}
See your doctor if:
\begin{itemize}
  \item Acute-onset nystagmus with vertigo, ataxia, dysarthria, or diplopia (suggests brainstem or cerebellar stroke)
  \item New nystagmus in adult without prior history
  \item Vertical nystagmus (especially downbeat or upbeat), which is highly localizing to CNS pathology
  \item Nystagmus with headache, neck stiffness, or altered mental status
  \item Nystagmus in a child with poor visual behavior or abnormal red reflex
\end{itemize}

\section*{Related Symptoms}
\begin{itemize}
  \item Vertigo
  \item Dizziness
  \item Oscillopsia (sensation that the visual world is bouncing)
  \item Imbalance
  \item Ataxia
  \item Diplopia
  \item Head tilt or turn
  \item Nausea and vomiting
\end{itemize}
\textbf{Important:} The direction and gaze-dependence of nystagmus provide critical localizing value -- downbeat nystagmus suggests craniocervical junction pathology (Chiari), while direction-changing gaze-evoked nystagmus implies cerebellar or brainstem integrator dysfunction.

\section*{Self-Care / Home Management}
\begin{itemize}
  \item Avoid driving or operating machinery during acute episodes, as oscillopsia impairs visual acuity.
  \item Establish a safe home environment with adequate lighting and clear pathways to prevent falls.
  \item Practice visual fixation on a near target to dampen nystagmus intensity during symptomatic periods.
  \item Limit alcohol, caffeine, and sedating medications which may worsen nystagmus.
\end{itemize}

\section*{How Your Doctor Will Evaluate You}
\begin{itemize}
  \item Observing eye movements in primary gaze, eccentric gaze positions, and during vergence identifies nystagmus type, direction, and characteristics.
  \\item Frenzel goggles eliminate visual fixation to reveal peripheral vestibular nystagmus.
  \item Neurologic examination focuses on cranial nerves, cerebellar function (dysmetria, dysdiadochokinesia), gait, and vestibular testing (head impulse test, Dix-Hallpike maneuver).
  \item MRI brain with and without contrast is essential for acquired nystagmus, focusing on posterior fossa and brainstem.
  \item Vestibular testing (VNG, ENG, caloric testing) differentiates peripheral from central etiologies.
\end{itemize}

\section*{Treatment Options}
\begin{itemize}
  \item Peripheral vestibular nystagmus: Treat underlying cause (canalith repositioning for BPPV, vestibular suppressants for acute neuritis, corticosteroids for M\'eni\`ere disease).
  \item Central nystagmus: Treat underlying causes (thiamine for Wernicke, immunosuppression for MS, surgical decompression for Chiari, stroke management).
  \item Pharmacologic suppression: Gabapentin, memantine, baclofen, clonazepam may reduce nystagmus amplitude and oscillopsia.
  \item Congenital nystagmus: Optical devices (contact lenses, prisms), visual acuity optimization, and experimental therapies (acetyl-leucine, tenotomy).
  \item Strabismus surgery (Kestenbaum-Anderson procedure) shifts the null zone to primary position in selected congenital cases.
\end{itemize}

\clearpage
